% Replacement hardware subsection, NOT included by the current manuscript.
% Setup/data facts below are verified. Protocol paragraphs are planned, not
% completed experiments. Replace TBD only from a frozen physical trial ledger.
% Keep the existing draft unchanged until these results are available.
\subsection{Real Robot Evaluation}
\label{sec:hardware}

\textbf{Setup and demonstrations.}
Our physical platform uses a 7-DoF xArm7 with a parallel gripper on a linear
rail. A second arm provides a stationary camera view. Both rails and the
camera arm remain fixed, while the policy observes the two wrist RGB cameras
and robot state. We collected handheld teleoperation demonstrations using
printed assembly fixtures. Lamp assembly requires grasping a hollow shade,
lowering it over a post, and releasing it in place. Drawer assembly requires
inserting and releasing two drawers in their designated slots.

The lamp dataset contains 52 demonstrations, with 47 used for training and
five held out. The drawer dataset contains 51 demonstrations, split 46/5.
We train a separate policy for each task. Both heads use the same pretrained
SmolVLA initialization, images, state, demonstrations and task caption. For
lamp assembly, each receives 40,000 updates with batch size 32. We retain
checkpoints every 10,000 updates and select each head's checkpoint using only
held-out command prediction errors around grasp closure and release.
% Confirm all runs completed before retaining the previous sentence as fact.
The spline head represents actions with eight cubic control points and an
8--24-frame event duration. Both heads return an eight-command execution
prefix in the native end-effector delta interface, with absolute gripper
opening targets.

\textbf{Planned evaluation.}
Before policy testing, a recorded demonstration will be replayed through the
execution interface, with its initial configuration and timing documented.
The learned policies will then be evaluated on matched starting layouts,
alternating the head tested first. Model weights and control settings will
be fixed before counted trials. Full lamp success requires the released shade
to remain seated for three seconds without human correction. Full drawer
success requires both drawers to be seated in their designated slots and
released. Stops and operator interventions count as failures for autonomous
completion. We will report grasp attainment separately from full success.
% Convert this paragraph to past tense only after the protocol is executed.
% Report actual executed cadence, checkpoint steps/hashes, trial counts, and
% any deviations. Do not imply that recorded 25 Hz equals the tested wire clock.

\begin{table}[t]
\centering
\caption{Real-robot evaluation template. No outcomes have been entered.
Each success entry will be reported as a count over all counted trials.}
\label{tab:hardware}
\setlength{\tabcolsep}{5pt}
\begin{tabular}{llcc}
\toprule
Task & Head & Grasp & Full success \\
\midrule
Lamp & Waypoint & TBD & TBD \\
     & Event spline & TBD & TBD \\
Drawer & Waypoint & TBD & TBD \\
       & Event spline & TBD & TBD \\
\bottomrule
\end{tabular}
\end{table}

% Figure: one setup crop identifying active arm / two views / printed fixture,
% plus a single successful policy filmstrip (approach, grasp, transfer, seat,
% release). Source every panel from a named counted trial; do not substitute a
% teleoperation demonstration for a learned rollout. Reserve ~0.35 page.
% If a matched physical-clock intervention finishes, replace optional drawer
% material with a small paired clock comparison. Use the same head checkpoint
% and matched waypoint interpolation/retiming baseline; report actual rates.
% Hardware currently uses whole-task captions. It tests physical deployment
% (and timing only if measured), NOT clause steering or novel program ordering.
